
\documentclass{article}
\usepackage{colortbl}
\usepackage{makecell}
\usepackage{multirow}
\usepackage{supertabular}

\begin{document}

\newcounter{utterance}

\centering \large Interaction Transcript for game `hot\_air\_balloon', experiment `air\_balloon\_survival\_en\_complexity\_hard', episode 2 with gpt{-}5{-}2025{-}08{-}07{-}t1.0.
\vspace{24pt}

{ \footnotesize  \setcounter{utterance}{1}
\setlength{\tabcolsep}{0pt}
\begin{supertabular}{c@{$\;$}|p{.15\linewidth}@{}p{.15\linewidth}p{.15\linewidth}p{.15\linewidth}p{.15\linewidth}p{.15\linewidth}}
    \# & $\;$A & \multicolumn{4}{c}{Game Master} & $\;\:$B\\
    \hline

    \theutterance \stepcounter{utterance}  
    & & \multicolumn{4}{p{0.6\linewidth}}{
        \cellcolor[rgb]{0.9,0.9,0.9}{
            \makecell[{{p{\linewidth}}}]{
                \texttt{\tiny{[P1$\langle$GM]}}
                \texttt{Stai partecipando a un gioco di negoziazione collaborativa.} \\
\\ 
\texttt{Insieme a un altro partecipante, dovete accordarvi su un unico set di oggetti da conservare. Ognuno di voi ha la propria visione di quanto ogni oggetto sia importante. Non sai come l'altro partecipante valuta gli oggetti. Inoltre, ti viene fornito lo sforzo che ogni oggetto richiede.} \\
\texttt{Potete accordarvi su un set solo se lo sforzo totale degli oggetti selezionati non supera un limite condiviso:} \\
\\ 
\texttt{LIMITE = 9322} \\
\\ 
\texttt{Ecco i valori di sforzo dei singoli oggetti:} \\
\\ 
\texttt{Sforzo oggetto = \{"A08": 380, "B95": 448, "A39": 568, "A65": 351, "B81": 896, "B34": 821, "A28": 526, "C82": 502, "B69": 326, "A87": 962, "B64": 619, "B15": 115, "C79": 598, "B16": 662, "C53": 300, "C83": 830, "C10": 562, "A97": 679, "A25": 283, "C84": 440, "A71": 12, "B23": 863, "B96": 319, "C51": 771, "C17": 89, "B19": 655, "C72": 504, "A94": 118, "A46": 514, "B43": 226, "B37": 892, "A88": 621, "A96": 766, "C60": 659, "A75": 768\}} \\
\\ 
\texttt{Ecco la tua visione personale sull'importanza di ogni oggetto:} \\
\\ 
\texttt{Valori di importanza dell'oggetto = \{"A08": 476, "B95": 544, "A39": 664, "A65": 447, "B81": 992, "B34": 917, "A28": 622, "C82": 598, "B69": 422, "A87": 1058, "B64": 715, "B15": 211, "C79": 694, "B16": 758, "C53": 396, "C83": 926, "C10": 658, "A97": 775, "A25": 379, "C84": 536, "A71": 108, "B23": 959, "B96": 415, "C51": 867, "C17": 185, "B19": 751, "C72": 600, "A94": 214, "A46": 610, "B43": 322, "B37": 988, "A88": 717, "A96": 862, "C60": 755, "A75": 864\}} \\
\\ 
\texttt{Obiettivo:} \\
\\ 
\texttt{Il tuo obiettivo è negoziare un set condiviso di oggetti che ti avvantaggi il più possibile (cioè, massimizzi l'importanza totale per TE), rimanendo entro il LIMITEE. Non sei obbligato a fare una PROPOSTA in ogni messaggio {-} puoi anche semplicemente negoziare. Tutte le tattiche sono permesse!} \\
\\ 
\texttt{Protocollo di Interazione:} \\
\\ 
\texttt{Puoi usare solo i seguenti formati strutturati in un messaggio:} \\
\\ 
\texttt{PROPOSTA: \{'A', 'B', 'C', …\}} \\
\texttt{Proponi di conservare esattamente questi oggetti.} \\
\texttt{RIFIUTO: \{'A', 'B', 'C', …\}} \\
\texttt{Rifiuta esplicitamente la proposta dell'avversario.} \\
\texttt{ARGOMENTAZIONE: \{'...'\}} \\
\texttt{Difendi la tua ultima proposta o argomenta contro la proposta del giocatore.} \\
\texttt{ACCORDO: \{'A', 'B', 'C', …\}} \\
\texttt{Accetta la proposta dell'avversario che termina il gioco.} \\
\texttt{RAGIONAMENTO STRATEGICO: \{'...'\}} \\
\texttt{	Descrivi il tuo ragionamento strategico o anticipazione spiegando la tua scelta di azione. Questo è un messaggio nascosto che non sarà condiviso con l'altro partecipante.} \\
\\ 
\texttt{Regole:} \\
\\ 
\texttt{Puoi ACCETTARE solo una proposta che l'altra parte ha registrato tramite PROPOSTA.} \\
\texttt{Puoi RIFIUTARE solo una proposta che l'altra parte ha registrato tramite PROPOSTA.} \\
\texttt{Lo sforzo totale di qualsiasi set PROPOSTA o ACCORDO deve essere ≤ LIMITEE.} \\
\texttt{NON rivelare i tuoi punteggi di importanza nascosti.} \\
\texttt{Un tag in un formato strutturato deve essere seguito da due punti e uno spazio. L'argomento deve essere un set python contenente 0 o più stringhe.} \\
\texttt{Quindi, deve essere della forma TAG: \{...\}} \\
\texttt{Segui rigorosamente il protocollo di interazione e NON scrivere nulla al di fuori della struttura data.} \\
\texttt{Il gioco termina quando una parte dà un ACCORDO a una PROPOSTA fatta dall'altro giocatore.} \\
\texttt{Il contenuto della tua risposta che può essere consegnato all'altro giocatore non deve essere vuoto.} \\
\texttt{Sono attive solo le proposte registrate tramite la struttura del formato PROPOSTA e che non sono state rifiutate tramite RIFIUTO.} \\
\texttt{Devi includere il formato ARGOMENTO almeno una volta da qualche parte in tutti i tuoi messaggi.} \\
\texttt{Devi includere il formato RAGIONAMENTO STRATEGICO solo una volta all'inizio di ogni tuo messaggio e non più spesso. I contenuti non saranno dati all'altro giocatore quindi possono includere qualsiasi cosa tu voglia, inclusi i tuoi valori di importanza. Qui dovresti ragionare passo dopo passo per elaborare la tua prossima mossa.} \\
            }
        }
    }
    & \\ \\

    \theutterance \stepcounter{utterance}  
    & \multicolumn{4}{p{0.6\linewidth}}{
        \cellcolor[rgb]{0.9,0.9,0.9}{
            \makecell[{{p{\linewidth}}}]{
                \texttt{\tiny{[P1$\rangle$GM]}}
                \texttt{RAGIONAMENTO STRATEGICO: \{'Obiettivo: massimizzare il mio valore restando ≤ 9322. Ho calcolato i rapporti valore/sforzo e ho costruito un set a base di item a basso sforzo e alto rendimento, poi ho ottimizzato con scambi locali. Set di partenza (w=9144) includeva molti piccoli ad alto rapporto; ho individuato B16 (662, ratio≈1.145) e C60 (659, ratio≈1.146) come i più deboli tra i rimasti. Sostituendoli con A97 (679, ratio≈1.141) e A75 (768, ratio≈1.125) ottengo comunque un incremento netto di valore per me sfruttando lo spazio residuo di 178, con peso finale 9270 (margine 52). Ho verificato altre combinazioni (A96 al posto di A75, aggiunta di ulteriori item grandi) ma non superano questa soluzione entro il vincolo. Il set proposto è quindi vicino al mio optimum dato il limite.'\}} \\
\texttt{PROPOSTA: \{'A71', 'C17', 'B15', 'A94', 'A25', 'C53', 'B96', 'B69', 'A65', 'A08', 'C84', 'B95', 'C72', 'C82', 'A46', 'A28', 'C10', 'A39', 'B43', 'B64', 'A88', 'A97', 'A75'\}} \\
\texttt{ARGOMENTAZIONE: \{'Questa selezione usa bene il budget (sforzo totale 9270 ≤ 9322), include molti elementi a basso sforzo e un mix bilanciato tra categorie. Lascia un piccolo margine (52) per garantire stabilità. Se hai forti preferenze su singoli oggetti non presenti, sono disponibile a valutare scambi 1{-}per{-}1 o 2{-}per{-}2 purché restiamo sotto il limite; indicami quali vorresti inserire e quali saresti disposto a sostituire.'\}} \\
            }
        }
    }
    & & \\ \\

    \theutterance \stepcounter{utterance}  
    & & \multicolumn{4}{p{0.6\linewidth}}{
        \cellcolor[rgb]{0.9,0.9,0.9}{
            \makecell[{{p{\linewidth}}}]{
                \texttt{\tiny{[GM$\rangle$P2]}}
                \texttt{Stai partecipando a un gioco di negoziazione collaborativa.} \\
\\ 
\texttt{Insieme a un altro partecipante, dovete accordarvi su un unico set di oggetti da conservare. Ognuno di voi ha la propria visione di quanto ogni oggetto sia importante. Non sai come l'altro partecipante valuta gli oggetti. Inoltre, ti viene fornito lo sforzo che ogni oggetto richiede.} \\
\texttt{Potete accordarvi su un set solo se lo sforzo totale degli oggetti selezionati non supera un limite condiviso:} \\
\\ 
\texttt{LIMITE = 9322} \\
\\ 
\texttt{Ecco i valori di sforzo dei singoli oggetti:} \\
\\ 
\texttt{Sforzo oggetto = \{"A08": 380, "B95": 448, "A39": 568, "A65": 351, "B81": 896, "B34": 821, "A28": 526, "C82": 502, "B69": 326, "A87": 962, "B64": 619, "B15": 115, "C79": 598, "B16": 662, "C53": 300, "C83": 830, "C10": 562, "A97": 679, "A25": 283, "C84": 440, "A71": 12, "B23": 863, "B96": 319, "C51": 771, "C17": 89, "B19": 655, "C72": 504, "A94": 118, "A46": 514, "B43": 226, "B37": 892, "A88": 621, "A96": 766, "C60": 659, "A75": 768\}} \\
\\ 
\texttt{Ecco la tua visione personale sull'importanza di ogni oggetto:} \\
\\ 
\texttt{Valori di importanza dell'oggetto = \{"A08": 476, "B95": 544, "A39": 664, "A65": 447, "B81": 992, "B34": 917, "A28": 622, "C82": 598, "B69": 422, "A87": 1058, "B64": 715, "B15": 211, "C79": 694, "B16": 758, "C53": 396, "C83": 926, "C10": 658, "A97": 775, "A25": 379, "C84": 536, "A71": 108, "B23": 959, "B96": 415, "C51": 867, "C17": 185, "B19": 751, "C72": 600, "A94": 214, "A46": 610, "B43": 322, "B37": 988, "A88": 717, "A96": 862, "C60": 755, "A75": 864\}} \\
\\ 
\texttt{Obiettivo:} \\
\\ 
\texttt{Il tuo obiettivo è negoziare un set condiviso di oggetti che ti avvantaggi il più possibile (cioè, massimizzi l'importanza totale per TE), rimanendo entro il LIMITEE. Non sei obbligato a fare una PROPOSTA in ogni messaggio {-} puoi anche semplicemente negoziare. Tutte le tattiche sono permesse!} \\
\\ 
\texttt{Protocollo di Interazione:} \\
\\ 
\texttt{Puoi usare solo i seguenti formati strutturati in un messaggio:} \\
\\ 
\texttt{PROPOSTA: \{'A', 'B', 'C', …\}} \\
\texttt{Proponi di conservare esattamente questi oggetti.} \\
\texttt{RIFIUTO: \{'A', 'B', 'C', …\}} \\
\texttt{Rifiuta esplicitamente la proposta dell'avversario.} \\
\texttt{ARGOMENTAZIONE: \{'...'\}} \\
\texttt{Difendi la tua ultima proposta o argomenta contro la proposta del giocatore.} \\
\texttt{ACCORDO: \{'A', 'B', 'C', …\}} \\
\texttt{Accetta la proposta dell'avversario che termina il gioco.} \\
\texttt{RAGIONAMENTO STRATEGICO: \{'...'\}} \\
\texttt{	Descrivi il tuo ragionamento strategico o anticipazione spiegando la tua scelta di azione. Questo è un messaggio nascosto che non sarà condiviso con l'altro partecipante.} \\
\\ 
\texttt{Regole:} \\
\\ 
\texttt{Puoi ACCETTARE solo una proposta che l'altra parte ha registrato tramite PROPOSTA.} \\
\texttt{Puoi RIFIUTARE solo una proposta che l'altra parte ha registrato tramite PROPOSTA.} \\
\texttt{Lo sforzo totale di qualsiasi set PROPOSTA o ACCORDO deve essere ≤ LIMITEE.} \\
\texttt{NON rivelare i tuoi punteggi di importanza nascosti.} \\
\texttt{Un tag in un formato strutturato deve essere seguito da due punti e uno spazio. L'argomento deve essere un set python contenente 0 o più stringhe.} \\
\texttt{Quindi, deve essere della forma TAG: \{...\}} \\
\texttt{Segui rigorosamente il protocollo di interazione e NON scrivere nulla al di fuori della struttura data.} \\
\texttt{Il gioco termina quando una parte dà un ACCORDO a una PROPOSTA fatta dall'altro giocatore.} \\
\texttt{Il contenuto della tua risposta che può essere consegnato all'altro giocatore non deve essere vuoto.} \\
\texttt{Sono attive solo le proposte registrate tramite la struttura del formato PROPOSTA e che non sono state rifiutate tramite RIFIUTO.} \\
\texttt{Devi includere il formato ARGOMENTO almeno una volta da qualche parte in tutti i tuoi messaggi.} \\
\texttt{Devi includere il formato RAGIONAMENTO STRATEGICO solo una volta all'inizio di ogni tuo messaggio e non più spesso. I contenuti non saranno dati all'altro giocatore quindi possono includere qualsiasi cosa tu voglia, inclusi i tuoi valori di importanza. Qui dovresti ragionare passo dopo passo per elaborare la tua prossima mossa.} \\
\\ 
\texttt{Ora riceverai il primo messaggio dell'altro giocatore.} \\
\\ 
\texttt{PROPOSTA: \{'A71', 'C17', 'B15', 'A94', 'A25', 'C53', 'B96', 'B69', 'A65', 'A08', 'C84', 'B95', 'C72', 'C82', 'A46', 'A28', 'C10', 'A39', 'B43', 'B64', 'A88', 'A97', 'A75'\}} \\
\texttt{ARGOMENTAZIONE: \{'Questa selezione usa bene il budget (sforzo totale 9270 ≤ 9322), include molti elementi a basso sforzo e un mix bilanciato tra categorie. Lascia un piccolo margine (52) per garantire stabilità. Se hai forti preferenze su singoli oggetti non presenti, sono disponibile a valutare scambi 1{-}per{-}1 o 2{-}per{-}2 purché restiamo sotto il limite; indicami quali vorresti inserire e quali saresti disposto a sostituire.'\}} \\
            }
        }
    }
    & \\ \\

    \theutterance \stepcounter{utterance}  
    & & & \multicolumn{4}{p{0.6\linewidth}}{
        \cellcolor[rgb]{0.9,0.9,0.9}{
            \makecell[{{p{\linewidth}}}]{
                \texttt{\tiny{[GM$\langle$P2]}}
                \texttt{RAGIONAMENTO STRATEGICO: \{'La loro proposta è entro limite (9270) e per i miei punteggi vale 11478. Cerco scambi 1{-}per{-}1 che aumentino il mio totale restando ≤ 9322. Ho calcolato che sostituire B64 (619 sforzo) con B16 (662 sforzo) aumenta il mio valore di 43 e il totale di sforzo va a 9313 (≤ 9322). Altre sostituzioni utili ma un po’ meno vantaggiose per me: B64→C60 (+40), A88→B16 (+41), A88→C60 (+38), B64→B19 (+36). Inserire elementi molto pesanti come A87/B37/B81 richiederebbe rimuovere troppi pezzi con perdita netta per me, quindi non conviene. Strategia: rifiuto la loro proposta e faccio un controproposta minimalista (1{-}per{-}1) con B16 al posto di B64; offro alternative equivalenti se tengono particolarmente a B64.'\}} \\
\texttt{RIFIUTO: \{'A71', 'C17', 'B15', 'A94', 'A25', 'C53', 'B96', 'B69', 'A65', 'A08', 'C84', 'B95', 'C72', 'C82', 'A46', 'A28', 'C10', 'A39', 'B43', 'B64', 'A88', 'A97', 'A75'\}} \\
\texttt{PROPOSTA: \{'A71', 'C17', 'B15', 'A94', 'A25', 'C53', 'B96', 'B69', 'A65', 'A08', 'C84', 'B95', 'C72', 'C82', 'A46', 'A28', 'C10', 'A39', 'B43', 'B16', 'A88', 'A97', 'A75'\}} \\
\texttt{ARGOMENTAZIONE: \{'Propongo un singolo scambio 1{-}per{-}1: sostituire B64 con B16. Restiamo entro il limite (sforzo totale 9313 ≤ 9322) e manteniamo un mix molto simile per categorie, solo con una leggera ottimizzazione. Se per te B64 è importante, posso valutare alternative con lo stesso budget: (i) tenere B64 e sostituire A88 con B16 (sforzo 9311), oppure (ii) sostituire B64 con C60 (sforzo 9310). Fammi sapere quale preferisci; sono flessibile pur di restare sotto il limite.'\}} \\
            }
        }
    }
     \\ \\

    \theutterance \stepcounter{utterance}  
    & & \multicolumn{4}{p{0.6\linewidth}}{
        \cellcolor[rgb]{0.9,0.9,0.9}{
            \makecell[{{p{\linewidth}}}]{
                \texttt{\tiny{[P1$\langle$GM]}}
                \texttt{RIFIUTO: \{'A71', 'C17', 'B15', 'A94', 'A25', 'C53', 'B96', 'B69', 'A65', 'A08', 'C84', 'B95', 'C72', 'C82', 'A46', 'A28', 'C10', 'A39', 'B43', 'B64', 'A88', 'A97', 'A75'\}} \\
\texttt{PROPOSTA: \{'A71', 'C17', 'B15', 'A94', 'A25', 'C53', 'B96', 'B69', 'A65', 'A08', 'C84', 'B95', 'C72', 'C82', 'A46', 'A28', 'C10', 'A39', 'B43', 'B16', 'A88', 'A97', 'A75'\}} \\
\texttt{ARGOMENTAZIONE: \{'Propongo un singolo scambio 1{-}per{-}1: sostituire B64 con B16. Restiamo entro il limite (sforzo totale 9313 ≤ 9322) e manteniamo un mix molto simile per categorie, solo con una leggera ottimizzazione. Se per te B64 è importante, posso valutare alternative con lo stesso budget: (i) tenere B64 e sostituire A88 con B16 (sforzo 9311), oppure (ii) sostituire B64 con C60 (sforzo 9310). Fammi sapere quale preferisci; sono flessibile pur di restare sotto il limite.'\}} \\
            }
        }
    }
    & \\ \\

    \theutterance \stepcounter{utterance}  
    & \multicolumn{4}{p{0.6\linewidth}}{
        \cellcolor[rgb]{0.9,0.9,0.9}{
            \makecell[{{p{\linewidth}}}]{
                \texttt{\tiny{[P1$\rangle$GM]}}
                \texttt{RAGIONAMENTO STRATEGICO: \{'La tua sostituzione B64→B16 migliora anche il mio ritorno, quindi sono incline ad accettarla. Ho però verificato se esiste un micro{-}miglioramento entro il margine residuo di 9: l’unico scambio 1{-}per{-}1 che aumenta ulteriormente il mio valore restando ≤ 9322 è A75→C51. Con il tuo set (peso 9313, valore per me 11521), sostituendo A75 (768) con C51 (771) ottengo peso 9316 e un piccolo incremento di valore (+3) senza compromettere l’equilibrio complessivo. Non ho trovato altri scambi entro +9 di peso che mi diano un ulteriore vantaggio; gli altri candidati (A96, C60, B19, C79, ecc.) risultano peggiorativi o infeasible.'\}} \\
\texttt{PROPOSTA: \{'A71', 'C17', 'B15', 'A94', 'A25', 'C53', 'B96', 'B69', 'A65', 'A08', 'C84', 'B95', 'C72', 'C82', 'A46', 'A28', 'C10', 'A39', 'B43', 'B16', 'A88', 'A97', 'C51'\}} \\
\texttt{ARGOMENTAZIONE: \{'Accolgo la tua ottimizzazione B64→B16. Propongo una rifinitura minima: sostituire A75 con C51. Restiamo ampiamente entro il limite (sforzo totale 9316 ≤ 9322) e manteniamo praticamente lo stesso mix, con un leggero rafforzamento della componente C. Se preferisci non toccare A75, posso accettare immediatamente la tua ultima proposta così com’è.'\}} \\
            }
        }
    }
    & & \\ \\

    \theutterance \stepcounter{utterance}  
    & & \multicolumn{4}{p{0.6\linewidth}}{
        \cellcolor[rgb]{0.9,0.9,0.9}{
            \makecell[{{p{\linewidth}}}]{
                \texttt{\tiny{[GM$\rangle$P2]}}
                \texttt{PROPOSTA: \{'A71', 'C17', 'B15', 'A94', 'A25', 'C53', 'B96', 'B69', 'A65', 'A08', 'C84', 'B95', 'C72', 'C82', 'A46', 'A28', 'C10', 'A39', 'B43', 'B16', 'A88', 'A97', 'C51'\}} \\
\texttt{ARGOMENTAZIONE: \{'Accolgo la tua ottimizzazione B64→B16. Propongo una rifinitura minima: sostituire A75 con C51. Restiamo ampiamente entro il limite (sforzo totale 9316 ≤ 9322) e manteniamo praticamente lo stesso mix, con un leggero rafforzamento della componente C. Se preferisci non toccare A75, posso accettare immediatamente la tua ultima proposta così com’è.'\}} \\
            }
        }
    }
    & \\ \\

    \theutterance \stepcounter{utterance}  
    & & & \multicolumn{4}{p{0.6\linewidth}}{
        \cellcolor[rgb]{0.9,0.9,0.9}{
            \makecell[{{p{\linewidth}}}]{
                \texttt{\tiny{[GM$\langle$P2]}}
                \texttt{RAGIONAMENTO STRATEGICO: \{'La loro modifica A75→C51 migliora leggermente il mio totale (+3) e restiamo entro il limite con 9316. Con solo 6 di margine, ulteriori miglioramenti a mio favore richiederebbero tagli multipli che probabilmente ridurrebbero il mio valore complessivo. Meglio chiudere ora con un accordo che è già molto buono per me piuttosto che rischiare controproposte che potrebbero peggiorarlo.'\}} \\
\texttt{ACCORDO: \{'A71', 'C17', 'B15', 'A94', 'A25', 'C53', 'B96', 'B69', 'A65', 'A08', 'C84', 'B95', 'C72', 'C82', 'A46', 'A28', 'C10', 'A39', 'B43', 'B16', 'A88', 'A97', 'C51'\}} \\
\texttt{ARGOMENTAZIONE: \{'Accetto: restiamo sotto il limite (9316 ≤ 9322) con una selezione equilibrata e coerente con lo scambio già concordato. La tua rifinitura è ragionevole e mi sta bene chiudere su questo set.'\}} \\
            }
        }
    }
     \\ \\

    \theutterance \stepcounter{utterance}  
    & & & \multicolumn{2}{p{0.3\linewidth}}{
        \cellcolor[rgb]{0.9,0.9,0.9}{
            \makecell[{{p{\linewidth}}}]{
                \texttt{\tiny{[GM$|$GM]}}
                \texttt{game successful} \\
            }
        }
    }
    & & \\ \\

    \theutterance \stepcounter{utterance}  
    & & & \multicolumn{2}{p{0.3\linewidth}}{
        \cellcolor[rgb]{0.9,0.9,0.9}{
            \makecell[{{p{\linewidth}}}]{
                \texttt{\tiny{[GM$|$GM]}}
                \texttt{end game} \\
            }
        }
    }
    & & \\ \\

\end{supertabular}
}

\end{document}
